AlphaFold2 a résolu en grande partie le problème du repliement des protéines isolées, en exploitant les signaux de coévolution d'un alignement multiple de séquences. Mais il ne voit que des protéines : ni ADN, ni ARN, ni ions, ni ligands. AlphaFold3 (2024) lève cette limite en prédisant des complexes de toutes ces molécules, grâce à une architecture remaniée (Pairformer, module de diffusion sur les atomes). Ce projet explique ce que faisait AlphaFold2, pourquoi il est aujourd'hui dépassé, et comment utiliser AlphaFold3. Un déploiement d'AlphaFold2 (ColabFold~1.5.5 dans Docker, sur CPU) sert de référence : la calmoduline humaine (149 résidus) y est prédite avec deux lobes fiables (pLDDT moyen 82,8) dont l'orientation relative reste incertaine (PAE d'environ 20~Å), et sans les ions Ca$^{2+}$ qu'elle fixe. AlphaFold3, utilisé via AlphaFold Server faute de GPU, prédit la même protéine avec une confiance nettement plus élevée (pLDDT moyen 89,7, pTM 0,61 contre 0,49), en quelques minutes au lieu de deux heures. Mais l'orientation des deux lobes reste incertaine (PAE d'environ 14~Å), ce qui reflète la flexibilité réelle de la protéine ; une partie du gain vient aussi des gabarits utilisés par le serveur.